\chapter{Interpretation Pitfalls}

\label{ch:ch35_pitfalls}

\section{Cognitive Pitfalls}
Common errors include anchoring on the referral diagnosis, search satisfaction, premature closure, confirmation bias, overcalling incidental findings, and failing to compare prior studies. Fatigue, interruptions, and poor communication increase risk.

\subsection{High-Risk Scenarios for Cognitive Error}
\begin{itemize}
    \item \textbf{After-hours and overnight reads} -- fatigue reduces diagnostic accuracy by up to 20\%.
    \item \textbf{High-volume lists} -- read speed increases while attention to detail decreases.
    \item \textbf{Emergency studies} -- time pressure may cause incomplete search patterns.
    \item \textbf{Prior diagnosis provided} -- the referring clinician's differential may anchor the reader.
    \item \textbf{Multiple incidental findings} -- the reader may focus on one and neglect others.
\end{itemize}

\begin{pearlbox}[Clinical Pearl: The 10-second pause]
Before finalising any report, pause for 10 seconds and ask: ``Have I checked the edges, the bones, the soft tissues, and the priors?'' This simple habit catches a significant proportion of missed findings.
\end{pearlbox}

\section{Common Specific Pitfalls}

\subsection{Missed Pneumothorax}
The most commonly missed finding on chest radiography. Always check the apices, the costophrenic angles, and the deep sulcus on supine films. On CT, a small pneumothorax may be visible only on lung windows at the lung apex.

\subsection{Overlooked Free Air}
Free intraperitoneal air on a supine abdominal radiograph may be subtle. Look for the Rigler sign (air on both sides of the bowel wall), the football sign (large volume of air outlining the falciform ligament), or the continuous diaphragm sign. When clinical suspicion is high, CT is far more sensitive.

\subsection{Normal Appendix Called Abnormal}
A fluid-filled, non-compressible structure $>6$~mm in diameter does not always indicate appendicitis. Periappendiceal fat stranding, wall thickening, and an appendicolith increase the positive predictive value. Clinical correlation is essential.

\subsection{Pseudotumour of the Kidney}
A prominent column of Bertin (renal cortical tissue extending between medullary pyramids) can mimic a renal mass on ultrasound. On CT, it enhances identically to renal cortex on all phases and has no mass effect.

\subsection{Pseudofracture (Looser Zone)}
In osteomalacia, radiolucent bands perpendicular to the cortex may simulate fractures. These ``pseudofractures'' are typically bilateral and symmetrical, found in the pubic ribia, femoral neck, and scapula.

\subsection{Pitfalls in Breast Imaging}
Dense breast tissue reduces mammographic sensitivity. Architectural distortion may be the only sign of invasive lobular carcinoma. Always correlate with ultrasound and consider supplemental MRI in high-risk patients with dense breasts.

\subsection{Pitfalls in CT Colonography}
Residual stool can mimic polyps. Use tagging agents and prone/supine comparisons to differentiate. A true polyp maintains its position and shape between prone and supine acquisitions.

\begin{pearlbox}[Clinical Pearl: Prior imaging is your best friend]
Comparing with the most recent prior study eliminates the majority of ``incidentalomas'' and provides context for subtle findings. Always request priors before reading a new study.
\end{pearlbox}

\section{Technical Pitfalls}
Check patient identity, laterality, coverage, motion, phase, contrast timing, sequence selection, and reconstruction. A technically limited examination should be described as such.

\subsection{Common Technical Errors}
\begin{table}[h]
\centering
\caption{Technical pitfalls and their impact on diagnosis.}
\label{tab:technical_pitfalls}
\begin{tabular}{p{3.5cm} p{5cm} p{5cm}}
\toprule
\textbf{Technical Error} & \textbf{Consequence} & \textbf{Prevention} \\
\midrule
Inadequate coverage & Missed findings at the edge of the scan range & Verify scan range covers the entire region of interest \\
Patient motion & Blurred images that may obscure or simulate pathology & Use motion reduction strategies; repeat if severely degraded \\
Wrong contrast phase & Hepatocellular carcinoma may be invisible on portal venous phase; requires arterial phase & Protocol correctly for the clinical question; use bolus tracking \\
Wrong laterality labeled & Catastrophic misdiagnosis or wrong-site procedure & Verify laterality on every image and in every report \\
Wrong sequence selected & Edema invisible on T1; restricted diffusion not assessed & Ensure appropriate sequence protocol for the clinical question \\
\bottomrule
\end{tabular}
\end{table}

\section{Error Reduction}
Use structured search patterns, checklists for high-risk studies, peer discussion for uncertainty, and documented critical-result communication. Learning systems should examine near misses and system factors without relying solely on individual blame.

\subsection{Checklist for Common Missed Findings}
\begin{enumerate}
    \item \textbf{Chest X-ray}: Apices, costophrenic sulci, hila, heart borders, mediastinum, bones, soft tissues.
    \item \textbf{CT head}: Sulci, ventricles, midline shift, extra-axial collections, orbits, sinuses, bone.
    \item \textbf{CT abdomen}: Liver margins, adrenal glands, kidneys (both poles), free fluid, bowel wall thickening, mesentery.
    \item \textbf{CT pulmonary angiography}: Central and segmental arteries, heart chambers, mediastinum, chest wall.
    \item \textbf{Musculoskeletal}: Align all cortices, trace every bone, check joint spaces, assess soft tissues.
\end{enumerate}

\begin{pearlbox}[Impression:]
Interpretation pitfalls are avoidable. A systematic search pattern, comparison with prior imaging, peer review for uncertain cases, and awareness of cognitive biases will substantially reduce missed diagnoses and preventable harm.
\end{pearlbox}

\clearpage
